\chapter{基本局部等价}\label{ch:15}
% 原PDF第71页，第15章开始。
本章解释整体局部化论证背后的局部输入\cite{FG05,FG07,ABC24a}。其基本思想是：临界水平的球面 Kac–Moody 表示，等同于非分歧，也就是单值性平凡的 oper 空间上的层。这将\chref{14}的真空生成图景提升到整个 Kazhdan–Lusztig 范畴。

\section{从 Beilinson–Drinfeld 到局部等价}\label{sec:15.1}\index{基本局部等价|hyperpage}
\srcpage{72}
回顾证明 $\Aut_\sigma$ 的 Hecke 特征性质时使用的两个输入。第一，对点 $x\in X$，Beilinson–Drinfeld 局部化函子给出
\[\operatorname{Loc}_x:\widehat{\mathfrak g}_{\mathrm{crit},x}\text{-}\mathrm{Mod}^{L^+G_x}\longrightarrow\DMod_{\mathrm{crit}}(\Bun_G),\]
而中心通过将局部 oper 上的函数拉回到全局 oper 来作用：
\origdiagram{p072-square.pdf}{0.4850}
这里，上标 $\mathbb V$-生成表示由真空模生成的子范畴。第二，对 $V\in\Rep(\checkG)$，几何 Satake 与临界全局截面给出
\[\operatorname{Sat}_x(V)\star\mathbb V_{\mathrm{crit},x}\simeq F_x(V)\otimes_{\mathfrak z_{\mathrm{crit}}(D_x)}\mathbb V_{\mathrm{crit},x},\]
其中 $F_x(V)$ 是 $\Op_{\checkG}(D_x)\simeq\Spec\mathfrak z_{\mathrm{crit}}(D_x)$ 上与普遍 $\checkG$-局部系统关联的向量丛。基本局部等价用整个球面临界范畴的等价，替换这一真空生成的陈述。

\section{非分歧 oper}\label{sec:15.2}\index{非分歧 oper|hyperpage}\index{Kazhdan–Lusztig 范畴|hyperpage}\index{临界 Weyl 对象|hyperpage}
令 $D=\Spec k[[t]]$、$D^\times=\Spec k((t))$、$L^+G=G[[t]]$、$LG=G((t))$。记
\[\operatorname{KL}_{G,\mathrm{crit}}:=\widehat{\mathfrak g}_{\mathrm{crit}}\text{-}\mathrm{Mod}^{L^+G}\]
为临界 Kazhdan–Lusztig 范畴。FLE 中的谱空间是比 $D$ 上正则 oper 空间更大的纤维积

\begin{transnote}{易错预警}
非分歧只要求底层局部系统延拓到 $D$，并不要求 oper 约化在 $D$ 上也正则延拓；所以谱空间取纤维积 $\Op_{\checkG}(D^\times)\times_{\LocSys_{\checkG}(D^\times)}\LocSys_{\checkG}(D)$，即「在 $D^\times$ 上是 oper、在 $D$ 上只是局部系统」。它比 $\Op_{\checkG}(D)$ 大：后者的点还要在 $t=0$ 处满足一般位置条件。缺这一对比，就看不懂 15.2 为何要用 $\check\lambda$ 分层去度量「缺陷」。[依据：原书 p.72 本段；并见 [FG07]、[ABC24a]。]
\end{transnote}
\begin{definition*}{}
\[\Op^{\mathrm{unr}}_{\checkG}:=\Op_{\checkG}(D^\times)\times_{\LocSys_{\checkG}(D^\times)}\LocSys_{\checkG}(D).\]
\end{definition*}
因此，它的一个点是穿孔圆盘上的 oper，其底层 $\checkG$-局部系统可延拓到 $D$。映射 $\Op_{\checkG}(D)\to\Op^{\mathrm{unr}}_{\checkG}$ 不是同构：局部系统可以延拓，并不迫使 oper 约化也正则地延拓。

具体地，选择平凡化后，$D^\times$ 上的 oper 具有形式 $\nabla=d+\omega$，其中 $\omega\in\check{\mathfrak g}((t))\,dt$。非分歧意味着，经过某个 $\checkG((t))$-规范变换后，它成为 $d+\eta$，其中 $\eta\in\check{\mathfrak g}[[t]]\,dt$。但同一个规范变换未必使 $\eta\bmod\check{\mathfrak b}[[t]]$ 在 $t=0$ 处满足一般位置条件。缺陷由优势权 $\check\lambda\in X^*(\check T)^+$ 度量，等价地，由 $G$ 的优势余权度量。在相应的正则分层上，可以写为
\[\eta\bmod\check{\mathfrak b}[[t]]=\sum_{i\in I}t^{\langle\check\lambda,\check\alpha_i\rangle}\psi_i(t)f_i\,dt,\]
其中 $\psi_i\in k[[t]]$ 且 $\psi_i(0)\ne0$，$\alpha_i$ 是 $\check{\mathfrak g}$ 的简单根，$f_i\in\check{\mathfrak g}_{-\alpha_i}$ 是简单负根向量。将此轨迹记为 $\Op^{\mathrm{reg},\check\lambda}_{\checkG}$。在集合意义下，
\[\Op^{\mathrm{unr}}_{\checkG}(k)=\coprod_{\check\lambda\in X^*(\check T)^+}\Op^{\mathrm{reg},\check\lambda}_{\checkG}(k).\]
更准确地，$\Op^{\mathrm{unr}}_{\checkG}$ 具有分支或形式片 $\Op^{\mathrm{unr},\check\lambda}_{\checkG}$，且
\[\Op^{\mathrm{reg},\check\lambda}_{\checkG}\subset(\Op^{\mathrm{unr},\check\lambda}_{\checkG})_{\mathrm{red}}.\]
\srcpage{73}
\begin{theorem}{临界 FLE，点态形式\cite{FG07,ABC24a}}{15.2.1}
存在 DG 范畴的等价
\[\operatorname{FLE}_{G,\mathrm{crit}}:\widehat{\mathfrak g}_{\mathrm{crit}}\text{-}\mathrm{Mod}^{L^+G}\xrightarrow{\sim}\IndCoh(\Op^{\mathrm{unr}}_{\checkG}),\]
它对标准的 Frenkel–Gaitsgory 阿贝尔 $t$-结构是 $t$-正合的，这里省略惯常的 $*$ 型或 $!$ 型 $\IndCoh$ 选择。此等价与球面范畴的几何 Satake 作用相容。
\end{theorem}

这一相容性在 Weyl 对象上可以看出。令 $\mathbb W^{\check\lambda}_{\mathrm{crit}}$ 为相应最高权的临界 Weyl 对象，并令 $i_{\check\lambda}:\Op^{\mathrm{reg},\check\lambda}_{\checkG}\hookrightarrow\Op^{\mathrm{unr}}_{\checkG}$。Drinfeld–Sokolov 约化给出
\[\Psi(\mathbb W^{\check\lambda}_{\mathrm{crit}})\simeq\mathfrak z^{\mathrm{reg},\check\lambda}\simeq\End(\mathbb W^{\check\lambda}_{\mathrm{crit}}),\qquad
\Spec\mathfrak z^{\mathrm{reg},\check\lambda}\simeq\Op^{\mathrm{reg},\check\lambda}_{\checkG}.\]
在 FLE 下，这说明，在相应 $\IndCoh$ 约定下，$\mathbb W^{\check\lambda}_{\mathrm{crit}}$ 对应于 $(i_{\check\lambda})_*^{\IndCoh}\cO_{\Op^{\mathrm{reg},\check\lambda}_{\checkG}}$。因此，球面临界范畴的支撑恰好是非分歧 oper 轨迹；自守侧的非分歧对应于谱侧单值性平凡。

\section{Drinfeld–Sokolov 约化与 Whittaker 形式}\label{sec:15.3}\index{Drinfeld–Sokolov 约化|hyperpage}\index{半无限上同调|hyperpage}\index{BRST 约化|hyperpage}
FLE 的底层函子是 Drinfeld–Sokolov 约化\cite{FG05,Fre02}。选定标准的非退化加性特征标 $\chi:\mathfrak n((t))\to k$。对临界模 $M$，置
\[\Psi(M):=C^{\frac\infty2}(\mathfrak n((t)),M\otimes k_\chi).\]

\begin{transnote}{含糊点拨}
$\infty/2$ 指半无限维数：先取 $\mathfrak n((t))$ 在临界中心扩张下的真空模，再作半无限 Chevalley 上链，$k_\chi$ 是加性特征标给出的一维模。临界水平时 $\Psi$ 恰好等同于 Drinfeld–Sokolov 约化，因此 $\Psi(M)$ 可以当作 $M$ 的半无限不变量；缺掉这层，$\Psi$ 就只是一个记号。[来源：nLab "Drinfeld-Sokolov reduction"，https://ncatlab.org/nlab/show/Drinfeld-Sokolov+reduction（Last revised 2013-06-10；2026-11-17 访问）；并见原书 p.73 本式与 [FG05]。]
\end{transnote}
这就是 $M$ 的半无限上同调，也称 BRST 约化。在临界水平，它有一个增强版本，取值于穿孔圆盘上 oper 空间上的层。相应的 Whittaker 陈述为
\[\widehat{\mathfrak g}_{\mathrm{crit}}\text{-}\mathrm{Mod}^{\mathrm{Whit}}\simeq\IndCoh(\Op_{\checkG}(D^\times)).\]
从 Whittaker 形式转到球面／Kazhdan–Lusztig 子范畴，在谱侧相当于施加底层局部系统可延拓到 $D$ 的条件，因此将 $\Op(D^\times)$ 替换为 $\Op^{\mathrm{unr}}$。

Satake 作用有平行的谱描述。临界球面 Hecke 范畴被识别为局部系统 Steinberg 叠上的 ind-凝聚层：
\[\begin{split}
\DMod(L^+G\backslash LG/L^+G)_{\mathrm{crit}}\simeq{}\\
\IndCoh\bigl(\LocSys_{\checkG}(D)\times_{\LocSys_{\checkG}(D^\times)}\LocSys_{\checkG}(D)\bigr).
\end{split}\]
\enlargethispage{\baselineskip}
在 FLE 下，这个范畴的卷积变成 $\IndCoh(\Op^{\mathrm{unr}}_{\checkG})$ 上自然的对应作用。在 Frenkel–Gaitsgory 最初的证明中，这出现在阿贝尔范畴和阿贝尔几何 Satake 的层面；整体 Langlands 需要 DG／因子化表述。

\section{与量子局部 Langlands 的关系}\label{sec:15.4}\index{量子局部 Langlands|hyperpage}
临界 FLE 是量子局部几何 Langlands 的一个边界值\cite{FG05}。离开边界后，对偶水平满足关系
\[\kappa=\kappa_{G,\mathrm{crit}}+c\,\kappa_{\mathrm{Kil},G},\qquad
\check\kappa=\kappa_{\checkG,\mathrm{crit}}+c^{-1}\kappa_{\mathrm{Kil},\checkG}.\]

\begin{transnote}{含糊点拨}
对偶侧取 $c^{-1}$ 是量子 Langlands 的水平配方：$\kappa$ 与 $\check\kappa$ 配对成临界值。$c\to0$ 时 $\check G$-Whittaker 侧退化，大水平极限给出 Satake 对应，所以「量子局部 Langlands 的极限等价就是 FLE」；这里 $c$ 不是自由参数，而是两侧水平之和偏离临界值的量。[依据：原书 p.73 本式与随后一段；并见 [FG05]。]
\end{transnote}
当 $c\ne0$ 时，预期的局部等价具有形式
\[\widehat{\mathfrak g}_{\kappa}\text{-}\mathrm{Mod}^{L^+G}\simeq\Whit(D_{\check\kappa}(\Gr_{\checkG})).\]
它在边界上的对应形式是标准的 Satake 与 Casselman–Shalika 等价。在大水平极限中，球面与 Whittaker Grassmann 簇范畴退化为
\[D(\Gr_G)^{L^+G}_{\mathrm{sph}}\simeq\Rep(\checkG),\qquad\Whit(D(\Gr_{\checkG}))\simeq\Rep(G).\]
\srcpage{74}
当 $c\to0$ 时，$G$ 侧变为临界水平，而 $\checkG$-Whittaker 侧变为 oper 的经典谱几何；极限等价就是 FLE
\[\widehat{\mathfrak g}_{\mathrm{crit}}\text{-}\mathrm{Mod}^{L^+G}\simeq\IndCoh(\Op^{\mathrm{unr}}_{\checkG}).\]

这是范围更广的量子局部对应表中的一行：
\[\begin{aligned}
LG\text{-}\mathrm{Mod}_{\kappa}&\longleftrightarrow L\checkG\text{-}\mathrm{Mod}_{\check\kappa},\\
\widehat{\mathfrak g}_{\kappa}\text{-}\mathrm{Mod}&\longleftrightarrow\Whit(D_{\check\kappa}(L\checkG)),\\
\Whit(D_{\kappa}(LG))&\longleftrightarrow\widehat{\check{\mathfrak g}}_{\check\kappa}\text{-}\mathrm{Mod},\\
D_{\kappa}(\Gr_G)&\longleftrightarrow D_{\check\kappa}(\Gr_{\checkG}),\\
D_{\kappa}(\Fl_G)&\longleftrightarrow D_{\check\kappa}(\Fl_{\checkG}).
\end{aligned}\]
在 $\kappa=\kappa_{\mathrm{crit}}$ 时，预期整个局部 $LG$-模理论是 $\LocSys_{\checkG}(D^\times)$ 上的范畴层：
\[LG\text{-}\mathrm{Mod}_{\mathrm{crit}}\rightsquigarrow\operatorname{ShvCat}(\LocSys_{\checkG}(D^\times)).\]
因此，Kac–Moody 范畴是 oper 轨迹上的范畴：它的临界中心是 $\Op_{\checkG}(D^\times)$ 上的函数。经过 Whittaker／Drinfeld–Sokolov 约化后，这个 oper 上的范畴被识别为
\[\IndCoh(\Op_{\checkG}(D^\times)),\]
这与 Feigin–Frenkel 中心的识别 $\operatorname{Spf}Z_{\mathrm{crit}}\simeq\Op_{\checkG}(D^\times)$ 相容。

\section{因子化 FLE 及其与整体局部化的相容性}\label{sec:15.5}\index{基本局部等价!因子化形式|hyperpage}\index{谱 Poincaré 函子|hyperpage}
\thref{15.2.1}中的点态等价有 Ran 空间上的因子化版本\cite{ABC24a}：
\[\operatorname{FLE}_{G,\mathrm{crit},\mathrm{Ran}}:\operatorname{KL}_{G,\mathrm{crit},\mathrm{Ran}}\xrightarrow{\sim}\IndCoh(\Op^{\mathrm{unr}}_{\checkG})_{\mathrm{Ran}}.\]
它与到 $\Bun_G$ 的局部化函子及全局谱函子相容。基本图表为
\origdiagram{p074-square.pdf}{0.5419}
这里 $\mathbb L$ 表示整体 Langlands 函子。函子 $\operatorname{Poinc}^{\mathrm{spec}}$ 是沿如下对应的推前—拉回变换：这一对应比较全局局部系统与它在形式圆盘上的限制。对点 $x$，记
\[\Op_{\checkG}(X\setminus\{x\})^{\mathrm{unr}}:=\Op_{\checkG}(X\setminus\{x\})\times_{\LocSys_{\checkG}(X\setminus\{x\})}\LocSys_{\checkG}(X),\]
相应的谱对应为
\[\Op_{\checkG}(X\setminus\{x\})^{\mathrm{unr}}\longrightarrow\Op_{\checkG}(D_x^\times)\times_{\LocSys_{\checkG}(D_x^\times)}\LocSys_{\checkG}(D_x).\]
\srcpage{75}
后面证明 $\mathbb L$ 为等价时，将使用这个图表，以及系数函子和无穷水平 FLE：
\origdiagram{p075-square.pdf}{0.5464}

\section{多点情形与生成性}\label{sec:15.6}\index{局部化!多点情形|hyperpage}
对有限点集 $\mathbf x=\{x_1,\ldots,x_n\}$，因子化给出
\[\operatorname{KL}_{G,\mathrm{crit},\mathbf x}\simeq\bigotimes_{i=1}^n\widehat{\mathfrak g}_{\mathrm{crit},x_i}\text{-}\mathrm{Mod}^{L^+G_{x_i}}.\]
通常的局部化函子
\[\operatorname{Loc}_{\mathbf x}:\operatorname{KL}_{G,\mathrm{crit},\mathbf x}\longrightarrow\DMod_{\mathrm{crit}}(\Bun_G)\]
可以通过在穿孔曲线上加入全局谱参数来增强。更准确地，记 $\Op^{\mathrm{unr}}_{\checkG,\mathbf x}:=\prod_i\Op^{\mathrm{unr}}_{\checkG,x_i}$，增强函子示意地具有形式
\[\begin{split}
\operatorname{KL}_{G,\mathrm{crit},\mathbf x}\otimes_{\QCoh(\Op^{\mathrm{unr}}_{\checkG,\mathbf x})}
\QCoh(\Op_{\checkG}(X\setminus\mathbf x)^{\mathrm{unr}})\\
\longrightarrow\DMod_{\mathrm{crit}}(\Bun_G),
\end{split}\]
其中
\[\Op_{\checkG}(X\setminus\mathbf x)^{\mathrm{unr}}:=\Op_{\checkG}(X\setminus\mathbf x)\times_{\LocSys_{\checkG}(X\setminus\mathbf x)}\LocSys_{\checkG}(X).\]
等价地，$\operatorname{Rep}(\checkG)_{\mathbf x}$ 在 $\operatorname{Loc}_{\mathbf x}$ 的像上的作用经由 $\LocSys_{\checkG}(X)$ 分解。

当 $\mathbf x$ 足够大时，这些局部范畴生成尖点部分。一个有用的示意表述是：对每个拟紧开集 $U\subset\Bun_G$，函子
\[\operatorname{KL}_{G,\mathrm{crit},\mathrm{Ran}}\longrightarrow\DMod_{\mathrm{crit}}(\Bun_G)\longrightarrow\DMod_{\mathrm{crit}}(U)\]
的像，在限制到 $U$ 后生成尖点贡献。整体论证的其余部分是 Eisenstein 部分：通过归纳，以及 FLE 与 Eisenstein 级数和常数项函子的相容性，处理非尖点对象。
